\chapter{Arsitektur Data dan Aplikasi Target}

\section*{Tujuan Pembelajaran}
\begin{itemize}
  \item Merancang arsitektur data target yang menjamin satu identitas dan satu rekam medis untuk tiap pasien.
  \item Menyusun lanskap aplikasi target beserta keputusan rasionalisasi untuk tiap aplikasi.
  \item Mengevaluasi pilihan membangun sendiri, membeli, atau berlangganan SaaS (\textit{Software as a Service}) dengan kriteria yang eksplisit.
\end{itemize}

\section{Pendahuluan}

% TODO
% Cakupan: satu pasien tercatat dengan banyak nomor rekam medis, sehingga riwayat klinis terpecah
%   dan laporan direksi tidak dapat dipercaya.

\section{Identitas Pasien Tunggal dan Data Induk}

% TODO
% Cakupan: MPI (Master Patient Index); pola manajemen data induk pada tingkat konsep (registri,
%   konsolidasi, terpusat); aturan kepemilikan dan pembaruan data induk; prinsip satu sumber
%   kebenaran dari C1.
% Kasus Medikara: penyatuan identitas pasien dari lima rumah sakit dan 18 klinik.
% Pendalaman teknis: aturan pencocokan identitas deterministik dan probabilistik beserta risiko
%   salah gabung dan salah pisah.

\section{Arsitektur Data untuk Analitik dan Pelaporan}

% TODO
% Cakupan: beda data operasional dan data analitik; gudang data dan platform data pada tingkat
%   konsep; mengganti laporan spreadsheet dengan sumber yang terkelola; kualitas data sebagai
%   syarat laporan terpadu.

\section{Lanskap Aplikasi Target dan Rasionalisasi}

% TODO
% Cakupan: keputusan per aplikasi (pertahankan, ganti, gabungkan, hentikan); aplikasi bersama dan
%   aplikasi lokal; pemeriksaan sebelum menghentikan aplikasi.
% Kasus Medikara: konsolidasi lima SIMRS dan posisi Medikara+ serta ERP.

\section{Integrasi Target pada Tingkat Konsep}

% TODO
% Cakupan: dari point-to-point ke hub, platform integrasi, dan API pada tingkat konsep; integrasi
%   eksternal dengan SATUSEHAT dan BPJS Kesehatan; batas tanggung jawab dengan mata kuliah IT40613
%   yang membahas rinci desain dan tata kelola API.
% Pendalaman teknis: memetakan sumber data Medikara ke resource HL7 FHIR.

\section{Keputusan Membangun, Membeli, atau Berlangganan}

% TODO
% Cakupan: kriteria berbobot (kesesuaian kapabilitas, biaya total, waktu, ketergantungan vendor,
%   kepatuhan); risiko pihak ketiga dari sisi arsitektur (dibahas rinci di mata kuliah IT40813);
%   mendokumentasikan alasan keputusan.
% Kasus Medikara: pelajaran dari insiden kebocoran data lewat vendor pada 2025.

\section{Aktivitas: Merancang Data dan Aplikasi Target}

% TODO
% Tujuan aktivitas: menghasilkan rancangan data dan aplikasi target yang dapat ditelusuri ke
%   kapabilitas dan prinsip.
% Peran dalam tim: analis teknologi memimpin rancangan data dan aplikasi target; analis bisnis
%   memeriksa keterkaitan tiap keputusan dengan kapabilitas dan prinsip; pemodel menggambar
%   application cooperation view target dan information structure view; koordinator menyelaraskan
%   rancangan integrasi dengan tim IT40613.
% Langkah: rancang identitas pasien tunggal dan data induk; rancang alur data untuk laporan;
%   susun tabel rasionalisasi aplikasi; bandingkan minimal dua opsi untuk satu keputusan besar
%   dengan kriteria berbobot; modelkan kondisi target.
% Luaran dan tonggak capstone: arsitektur data dan aplikasi target (C3).
% Berkas sources yang direncanakan: sources/chapter09/templat-tabel-rasionalisasi-aplikasi.csv,
%   sources/chapter09/templat-matriks-bangun-beli-langganan.csv,
%   sources/chapter09/contoh-aturan-identitas-pasien.md

\section{Ringkasan}

% TODO

\section{Latihan}

\begin{enumerate}
  \item Jelaskan tanpa istilah teknis mengapa MPI (\textit{Master Patient Index}) dibutuhkan sebelum laporan manajemen terpadu dapat dipercaya.
  \item Susun lima kriteria untuk memilih antara membangun sendiri dan membeli SIMRS (Sistem Informasi Manajemen Rumah Sakit) baru. Beri bobot dan jelaskan alasannya.
  \item Sebuah aplikasi dinilai \textit{Eliminate}, tetapi masih dipakai satu unit kecil untuk laporan ke regulator. Apa yang perlu diperiksa sebelum aplikasi itu dihentikan?
\end{enumerate}
